% =====================================================================
%  04-identidad.tex - Identidad, sesiones y cuentas
% =====================================================================

\section{Identidad, sesiones y cuentas}\label{sec:impl-identidad}

Esta sección reúne la autenticación de cuentas y su ciclo de vida: sesiones vigentes y segundo factor, recuperación de contraseña, alta por invitación, vínculo de certificado del Owner con su primer factor opcional y directorio de miembros con revocación durable. Las primitivas Argon2id y TOTP se describen con los servicios criptográficos (\cref{sec:impl-s7-argon2,sec:impl-s8-totp}); aquí se presenta su uso por la aplicación HTTP. Las vistas Qadra de estos recorridos aplican los mecanismos comunes de sesión, borradores y conciliación (\cref{sec:impl-http-interfaz}), y cada subsección precisa sólo lo propio de su dominio.

\subsection{Identidad y sesiones vigentes}\label{sec:impl-http-identidad}

El alta inicial solo funciona mientras no existen usuarios. PostgreSQL serializa esa condición y confirma el primer Owner junto con su evento de auditoría. El correo se normaliza y la contraseña se almacena como PHC Argon2id; el enrolamiento entrega una sola vez el secreto TOTP y ocho códigos de recuperación. La tabla \texttt{users}, creada por la migración 0001 y ampliada por la 0025 con la generación de autenticación (\cref{tab:anexo-migraciones}), conserva UUID, correo único, rol, estado activo, secreto TOTP cifrado, hashes de recuperación y revisión optimista. \texttt{Aes\allowbreak Gcm\allowbreak Secret\allowbreak Protector} liga el secreto al UUID del usuario mediante datos autenticados bajo la KEK del entorno.

Login verifica la contraseña y crea un desafío Redis de cinco minutos. Cada intento de segundo factor consume ese desafío mediante \texttt{GETDEL}; también lo consume un rechazo. TOTP añade un reclamo atómico contra repetición, mientras que recuperación confirma el consumo del código junto con su auditoría en PostgreSQL. Antes de emitir una sesión se comprueban nuevamente la cuenta activa, el principal completo y la generación de autenticación, cuyo avance revoca sesiones y desafíos anteriores (\cref{sec:impl-miembros}). Cinco contraseñas rechazadas bloquean durante 15 minutos la clave normalizada del correo; incremento y expiración del contador son atómicos.

Las sesiones utilizan tokens opacos de 256 bits y una clave Redis derivada por SHA-256, sin guardar el token en claro. El backend local conserva por defecto \texttt{absolute\_only}: un límite absoluto de 24 horas que ninguna actividad prolonga. La opción \texttt{-{}-session-idle-seconds} de \texttt{serve}, o \texttt{TT\_SESSION\_IDLE\_SECONDS}, habilita explícitamente un límite de inactividad positivo de hasta 86\,400 segundos; una configuración inválida impide el arranque. No se selecciona una duración de inactividad por omisión. La consulta de estado y las peticiones ordinarias no registran actividad.

Redis conserva una estructura hash versionada con identidad, política, creación, última actividad y plazos. Los scripts Lua usan el reloj del servidor y comprueban el esquema completo, la política, un TTL positivo y la fecha exacta de expiración. La identidad JSON permanece opaca para Lua, de modo que la generación de autenticación mantiene su valor entero. Una clave con otro tipo, versión o conjunto de campos, sin expiración o con una política distinta exige autenticar de nuevo; no se repara asignando tiempo adicional. La comprobación mediante \texttt{PEXPIRETIME} requiere Redis 7 o posterior, mientras que el contrato de respaldo exige separadamente Redis 7.4. El fundamento se conserva en \rutarepo{docs/adr/0064-explicit-session-activity.md}.

\begin{table}[H]
  \centering
  \caption{Contrato HTTP de identidad.}
  \label{tab:impl-http-identidad}
  \begin{tabularx}{\linewidth}{@{}L{6.2cm} X@{}}
    \toprule
    \textbf{Ruta} & \textbf{Responsabilidad} \\
    \midrule
    \texttt{POST /api/v1/auth/bootstrap} & Crea el primer Owner y entrega su enrolamiento; \texttt{201}. \\
    \texttt{POST /api/v1/auth/login} & Verifica contraseña y entrega desafío MFA. \\
    \texttt{POST /api/v1/auth/mfa/totp} & Consume desafío, verifica TOTP y crea sesión. \\
    \texttt{POST /api/v1/auth/mfa/recovery} & Consume desafío y recuperación; crea sesión. \\
    \texttt{GET /api/v1/auth/me} & Devuelve identidad y rol vigentes. \\
    \texttt{GET /api/v1/auth/session} & Consulta identidad, política y plazos sin renovar; \texttt{200}. \\
    \texttt{POST /api/v1/auth/activity} & Registra actividad explícita dentro de una sesión vigente; \texttt{200}. \\
    \texttt{POST /api/v1/auth/logout} & Revoca la sesión; \texttt{204}. \\
    \texttt{POST /api/v1/users} & Owner crea y enrola otra cuenta; \texttt{201}. \\
    \bottomrule
  \end{tabularx}
\end{table}

Ambos mecanismos MFA devuelven el bearer y los mismos metadatos públicos que la consulta de sesión: usuario, política, tiempo del servidor y vencimientos absoluto y de inactividad. Los instantes se expresan en milisegundos Unix; los campos de inactividad son \texttt{null} en el modo absoluto. El campo \texttt{expires\_in\_seconds} representa los segundos enteros restantes hasta el primer vencimiento respecto del tiempo del servidor devuelto. Las respuestas de estado y actividad no incluyen un bearer nuevo ni la generación interna de autenticación; sus campos precisos se documentan en \rutarepo{docs/http-api.md}.

Los cuatro roles pueden consultar y registrar actividad de su propia sesión mediante bearer vigente, cuerpo vacío y ausencia de parámetros de consulta. Ambas respuestas usan \texttt{Cache-Control: no-store}. Una sesión inválida recibe \texttt{401}; aportar cuerpo o parámetros recibe \texttt{400}. La actividad explícita solo desplaza el plazo de inactividad hasta el límite absoluto y no recrea claves vencidas o revocadas. La consulta de estado, \texttt{/auth/me} y las demás peticiones protegidas conservan los plazos.

Cada operación protegida consulta la sesión en Redis, contrasta cuenta activa, principal completo y generación con PostgreSQL, y confirma después que la misma sesión sigue vigente en Redis. Registrar actividad añade una comparación atómica final antes de modificar su plazo. La creación de usuarios revalida además al Owner al confirmar. PostgreSQL y Redis no comparten una transacción distribuida: un cambio de cuenta posterior a su lectura puede coincidir con la actividad, pero la generación modificada impide autorizaciones posteriores. Una operación admitida antes de vencer puede completar una transacción ya iniciada.

Si falla la auditoría después de crear un desafío o sesión, la aplicación intenta retirar la credencial y devuelve error sin token. Si también falla la limpieza, la clave puede permanecer hasta su expiración. Logout revoca primero y no restaura la sesión si falla el evento posterior. En el alta directa, si se pierde la respuesta de enrolamiento, su material MFA no puede recuperarse; el alta por invitaci\'on dispone, en cambio, de preparaci\'on y recibo recuperables (\cref{sec:impl-invitaciones}).

\subsection{Recuperación de contraseña con activación explícita}\label{sec:impl-recuperacion-password}

La recuperación realiza el flujo alterno de CU-03 y la parte correspondiente de RF-01 (\cref{tab:cu-03,tab:rf-acceso}): un enlace temporal por correo permite fijar otra contraseña (\cref{fig:impl-recuperacion-estados}). Es distinta del código de recuperación MFA, que exige conocer la contraseña, y sus rutas, formulario y correo están deshabilitados por omisión (ADR 0066).

\begin{figure}[!htb]
  \centering
  \begin{tikzpicture}[
  est/.style={cajaClara, text width=2.5cm, minimum height=1.15cm, rounded corners=5pt},
  ev/.style={etq, align=center}]

  % Estados principales
  \node[est] (s0) at (1.4,-1.9) {\textbf{Solicitada}\\ respuesta 202; no espera búsqueda, emisión ni envío};
  \node[est] (s1) at (5.65,-1.9) {\textbf{Emitida}\\ PostgreSQL guarda solo el SHA-256; el correo lleva el enlace};
  \node[est] (s2) at (9.9,-1.9) {\textbf{Consumida}\\ contraseña nueva; la capacidad ya no sirve};
  \node[cajaOsc, text width=2.9cm, minimum height=1.15cm] (fin) at (14.35,-1.9)
    {Acceso con la contraseña nueva y el MFA ya enrolado};

  \draw[flecha] (s0) -- node[ev, above=3pt]{cuenta\\ activa y\\ cupo libre} (s1);
  \draw[flecha] (s1) -- node[ev, above=3pt]{capacidad\\ vigente y\\ contraseña} (s2);
  \draw[flecha] (s2) -- node[ev, above=3pt]{sin\\ sesión} (fin);

  % Entrega incierta
  \node[paso, text width=3.4cm, anchor=south] (inc) at (5.65,-0.95)
    {Entrega incierta: la capacidad sigue vigente, sin reenvío automático};
  \draw[flechaSuave] (inc.south) -- (s1.north);

  % Salidas
  \node[cajaGris, text width=2.5cm, minimum height=1.0cm] (n0) at (1.4,-3.95)
    {\textbf{Sin emisión}: cuenta desconocida o inactiva, o cupo ocupado; mismo resultado, sin otra tarea};
  \node[cajaError, text width=1.9cm, minimum height=1.0cm] (e1) at (4.35,-3.95) {\textbf{Vencida}\\ (vigencia\\ configurada)};
  \node[cajaError, text width=2.3cm, minimum height=1.0cm] (e2) at (6.85,-3.95) {\textbf{Cancelada}: entrega rechazada (mismo ID y digest) o restauración};
  \draw[flecha] (s0) -- (n0);
  \draw[flechaError] (s1.south -| e1) -- (e1);
  \draw[flechaError] (s1.south -| e2) -- (e2);

  % Instalación privada
  \node[paso, text width=6.6cm, anchor=north west] at (9.0,0.05)
    {\textbf{Instalación privada}: Resend; vigencia de 900 s; un enlace pendiente por cuenta; una solicitud por correo cada 300 s; espera total de correo de 5000 ms};

  % Transacción de consumo
  \node[pasoOsc, text width=7.0cm, anchor=north] (tx) at (12.25,-2.85)
    {\textbf{Una transacción}, tras los bloqueos de auditoría y de fila:\\
     1. vuelve a comprobar vigencia y generación (la inspección previa a Argon2id no autoriza);\\
     2. sustituye el hash Argon2id de la contraseña;\\
     3. revisión +1 y generación de autenticación +1;\\
     4. consume la capacidad y agrega el evento.\\
     Si falla la auditoría, se revierte todo. Rol, pertenencias y MFA se conservan.};
  \draw[flechaSuave] (s2.south) -- (s2.south |- tx.north);
\end{tikzpicture}

  \caption{Ciclo de una capacidad de recuperación. La respuesta pública es igual haya o no emisión; el consumo avanza la generación (\cref{fig:impl-identidad-estado}).}
  \label{fig:impl-recuperacion-estados}
\end{figure}

\paragraph{Capacidad y transporte.} La capacidad es un valor aleatorio de 32 bytes del que se guarda solo una huella, con UUID propio por emisión. Un enlace nuevo no reemplaza otro vigente. La solicitud admite un correo acotado, y un único cupo cubre la entrada pendiente y su ejecución. El consumidor comparte con HTTP el presupuesto bloqueante, cancelar la espera no libera el permiso y, al detener el servidor, se descartan las tareas no iniciadas y se esperan las ya comenzadas. El emisor usa un destino HTTPS fijo de Resend, sin redirecciones, proxies de entorno ni reintentos, con la capacidad en el fragmento de una URL de origen aprobado.

\paragraph{Formulario.} Qadra retira el fragmento, pide confirmar la contraseña y no guarda el enlace; ante una respuesta incierta no repite el envío: se comprueba el acceso o se pide otro enlace, que descarta las respuestas tardías. Con sesión activa conserva los borradores (\cref{sec:impl-qadra-borradores}) antes de exigir otro acceso con MFA.

\paragraph{Configuración.} La activación exige habilitación expresa, remitente, origen, vigencia y cupo sin valor por omisión, cuotas, ventanas y tiempos del transporte; una clave de correo compartida no activa ningún canal, y los valores incompletos impiden abrir el servidor.

\paragraph{Límites.} Admitir una solicitud no es entrega durable ni prueba que la cuenta exista, y el resultado igual para cuentas desconocidas o inactivas no acredita uniformidad temporal del transporte. El rol operativo solo invoca operaciones restringidas, sin leer capacidades ni hashes. La habilitación operativa depende del transporte, de los límites y URL, de la retención y de la restauración: tras restaurar un respaldo, un comando administrativo (\cref{tab:impl-cli-catalogo}) que verifica la base, el esquema y la cabeza de auditoría esperados cancela las capacidades pendientes con un recibo, sin tocar credenciales, MFA ni capacidades consumidas. La verificación y la aceptación productiva están en la \cref{sec:pruebas-recuperacion-password}.

\Fuentes{\rutarepo{docs/password-reset-public.md}; ADR 0066.}


\subsection{Invitaciones y enrolamiento recuperable}\label{sec:impl-invitaciones}

La invitación realiza el alta de integrantes de CU-02 y RF-03 (\cref{tab:cu-02,tab:rf-acceso}; diseño y estados en las \cref{fig:cu-seq-alta,fig:impl-invitacion-estados}): el destinatario elige contraseña y MFA sin entregar esos secretos al Owner, que indica solo correo y perfil Litigante o Paralegal. Una cuenta existente, activa o inactiva, no se invita.

\begin{figure}[!htb]
  \centering
  \begin{tikzpicture}[
  est/.style={cajaClara, text width=3.0cm, minimum height=1.15cm, rounded corners=5pt},
  ent/.style={caja, text width=2.45cm, minimum height=0.6cm, rounded corners=5pt},
  ev/.style={etq, align=center}]

  % Estados no aceptados
  \draw[zona] (0.15,0.6) rectangle (8.65,-2.65);
  \node[zonaLabel] at (0.15,0.6) {NO ACEPTADAS};
  \node[est] (i0) at (2.0,-1.7) {\textbf{Pendiente}\\ capacidad de 32 B, revisión \emph{n}, vigencia de 24 h};
  \node[est] (i1) at (6.75,-1.7) {\textbf{Pendiente de MFA}\\ contraseña elegida; TOTP y códigos cifrados};
  \node[cajaOsc, text width=3.0cm, minimum height=1.15cm, rounded corners=5pt] (i2) at (11.6,-1.7)
    {\textbf{Aceptada}\\ cuenta, consumo, auditoría y recibo en una transacción};
  \draw[flecha] (i0) -- node[ev, above=3pt]{preparar} node[ev, below=3pt]{(contraseña)} (i1);
  \draw[flecha] (i1) -- node[ev, above=3pt]{finalizar\\ (TOTP)} (i2);
  \draw[flechaSuave] (i1.north west) .. controls +(0.1,0.85) and +(-0.1,0.85) .. node[etq, above, align=center]{misma contraseña y capacidad\\ vigente: mismo material} (i1.north east);

  % Salidas de cualquier estado no aceptado
  \node[cajaError, text width=2.6cm, minimum height=0.9cm] (i3) at (4.6,-3.55)
    {\textbf{Vencida}:\\ 24 h sin aceptar};
  \node[cajaGris, text width=2.6cm, minimum height=0.9cm] (i4) at (8.0,-3.55)
    {\textbf{Cancelada} tras restaurar una copia};
  \draw[flechaError] (4.6,-2.65) -- (i3);
  \draw[flecha, draw=ESCOMGris] (8.0,-2.65) -- node[etq, right=2pt]{restauración} (i4);

  % Reenvío
  \node[paso, text width=8.4cm, anchor=north, align=center] (re) at (5.0,-4.75)
    {\textbf{Reenvío explícito} (UUID nuevo y revisión esperada): mismo ID, correo y rol;\\ revisión \emph{n}+1, otra capacidad y otras 24 h; invalida el enlace y la preparación anteriores};
  \draw[flecha] (i3) -- (i3 |- re.north);
  \draw[flecha] (i4) -- (i4 |- re.north);
  \draw[flecha] (2.75,-2.65) -- (2.75,0 |- re.north);
  \draw[flecha] (1.0,0 |- re.north) -- (1.0,0 |- i0.south);

  % Entrega por revisión
  \draw[zona] (10.25,-2.75) rectangle (16,-6.15);
  \node[zonaLabel] at (10.25,-2.75) {ENTREGA DE CADA REVISIÓN};
  \node[ent, text width=1.9cm, minimum height=1.0cm] (d0) at (11.45,-4.45) {\textbf{Pendiente}\\ reservada hasta 30 s};
  \node[ent] (d1) at (14.5,-3.55) {\textbf{Enviada}: el proveedor aceptó};
  \node[ent] (d2) at (14.5,-4.45) {\textbf{Incierta}: otro intento tras 60 s};
  \node[ent, draw=IPNGuinda, text=IPNGuinda] (d3) at (14.5,-5.35) {\textbf{Rechazada}: exige reenvío};
  \foreach \d in {d1,d2,d3} \draw[flecha] (d0.east) -- (\d.west);
\end{tikzpicture}

  \caption{Estados de una invitación y de su entrega. Preparar no crea una cuenta utilizable; solo finalizar la crea, y después se entra con contraseña y MFA.}
  \label{fig:impl-invitacion-estados}
\end{figure}

\paragraph{Operaciones.} PostgreSQL guarda la huella de la capacidad (\cref{tab:anexo-canones}); el material del correo se cifra con AES-GCM ligado a invitación, revisión y finalidad, y el Owner solo consulta metadatos y estados. Preparación, confirmación y reenvío llevan UUID de operación distintos: una respuesta incierta se reconcilia repitiendo el comando exacto o con el recibo, que verifica la contraseña sin reentregar secretos MFA; cambiar el contenido de una operación aplicada da conflicto.

\paragraph{Cliente.} Qadra retira la capacidad del fragmento de la URL y solo persiste identificadores operativos y fase, nunca token, contraseña ni secretos MFA; el correo original permite retomar tras cerrar la página.

\paragraph{Límites.} La aceptación del proveedor no acredita la llegada al buzón. La activación es independiente, deshabilitada por omisión, y exige remitente, origen HTTPS y secreto del proveedor. Las rutas públicas, con límite propio de intentos, no crean ni adjuntan sesiones. Tras restaurar una copia, con admisión cerrada y escritores y consumidores detenidos, una operación administrativa expresa con la comprobación de la \cref{tab:impl-cli-catalogo}, nunca automática al arrancar, anula las invitaciones pendientes con su preparación MFA y sus entregas, sin tocar cuentas ni aceptaciones. Se verifica en la \cref{sec:pruebas-invitaciones}.

\Fuentes{\rutarepo{docs/member-invitations-api.md}; ADR 0075.}


\subsection{Vínculo explícito de certificado del Owner}\label{sec:impl-vinculo-certificado}

El vínculo relaciona una cuenta Owner consigo misma con un despliegue, una raíz, una revisión de confianza y un certificado público de la CA interna (\cref{fig:cripto-pki}). Sobre él se apoya el primer factor opcional por certificado de CU-03 (\cref{tab:cu-03}). El vínculo y el acceso usan declaraciones firmadas de longitud fija con propósitos separados (\cref{tab:anexo-canones}): \texttt{OWNCERT1}, de 150 bytes, registra o retira el vínculo con revisión de cuenta y generación; por sí solos, sus valores no autentican ni acreditan una escritura durable.

\paragraph{Verificador.} Bajo la sesión MFA y antes de los bloqueos, el perfil Partner exige RSA de 3072 bits, exponente 65537, SHA-256/RSA y usos de clave exactos con clientAuth y emailProtection, que el perfil de la declaración de participantes rechaza (\cref{sec:impl-declaracion-participante,sec:pki-declaraciones}). Recalcula el resumen SHA-256, verifica firma, raíz y CRL obligatoria al instante pedido y compara confianza y huellas: prueba consistencia con la confianza aportada, no su publicación ni su vigencia futura.

\paragraph{Persistencia y consulta.} PostgreSQL toma el bloqueo compartido de auditoría y el de la cuenta, relee Owner activo, contadores y confianza publicada y revalida la vigencia con el reloj leído tras las esperas. Alta y retiro solo añaden filas, sin modificar las anteriores, con su evento en la misma transacción; los disparadores no verifican RSA ni autentican. Cada cuenta tiene como máximo un vínculo sin retirar y una huella no pasa a otro Owner. El retiro es definitivo, conserva la captura y procede aunque el certificado haya vencido. La consulta del vínculo actual devuelve el recibo o indica que no hay vínculo, y falla si encuentra más de uno sin retirar; no audita ni verifica RSA o confianza vigente, así que su respuesta no impide un retiro posterior ni acredita vigencia.

\paragraph{Vista \emph{Mi certificado}.} Usa un UUID por intención y un único PEM público; firma con el ayudante local, sin que clave ni contraseña salgan del navegador, o importa una firma externa (\cref{sec:impl-criptografia-local}). Una escritura incierta se captura antes de esperar y se concilia por UUID exacto, sin repetirla salvo decisión expresa (\cref{sec:impl-qadra-conciliacion}); al reingresar solo queda en memoria el material público admitido (\cref{sec:impl-qadra-borradores}), y el recibo descargable no presenta la confianza histórica como vigente.

\paragraph{Acceso por certificado con MFA.} \texttt{OWNAUTH1} (\cref{tab:anexo-canones}) fija cuenta Owner, generación, vínculo y confianza; su intervalo de 1 a 300 s, con vencimiento exclusivo, solo limita presentar la prueba, y la validez de certificado y confianza es otra frontera. La captura en Redis vence en el instante firmado (\cref{fig:impl-acceso-certificado}). La revisión de cuenta puede avanzar al consumir un código de recuperación, pero generación y procedencia deben coincidir. La sesión guarda la revisión de confianza y sus plazos quedan acotados por la vigencia de la credencial; los registros de contraseña no cambian (ADR 0068).

\begin{figure}[!htb]
  \centering
  \begin{tikzpicture}[yscale=0.79,
  tramo/.style={draw=ESCOMGris, dotted, line width=0.5pt},
  tramoL/.style={font=\diag\bfseries, text=ESCOMAzulOscuro, fill=white,
    anchor=west, inner sep=1.5pt}]
  \Actores{usr/1.5/{Owner\\ (Qadra)},api/5.6/{API},rds/9.9/{Redis},pg/14.3/{PostgreSQL}}{-12.55}

  % ==== Declaración ====
  \draw[tramo] (0,-0.85) -- (16,-0.85);
  \node[tramoL] at (2.95,-0.85) {Declaración nueva};
  \Msg{1.5}{5.6}{-1.35}{\textbf{1} ¿acceso habilitado?}
  \Msg[msgR]{5.6}{1.5}{-1.8}{solo un booleano}
  \Paso{1.9}{-1.88}{3.5}{Elige el recibo público de su vínculo; conserva cuenta, vínculo, huella y nombre}
  \Msg{1.5}{5.6}{-3.3}{\textbf{2} pide declaración}
  \Msg{5.6}{9.9}{-3.75}{cuotas de inicio}
  \draw[msg] (5.6,-4.2) -- node[msgL, pos=0.74]{Owner activo, vínculo y confianza} (14.3,-4.2);
  \Msg{5.6}{9.9}{-4.65}{captura de un uso; vence al instante firmado}
  \Msg[msgR]{5.6}{1.5}{-5.1}{\texttt{OWNAUTH1} (182 B)}
  \Paso{1.9}{-5.25}{3.5}{Firma los bytes exactos con el ayudante local u OpenSSL}

  % ==== Prueba ====
  \draw[tramo] (0,-6.5) -- (16,-6.5);
  \node[tramoL] at (2.95,-6.5) {Prueba y segundo factor};
  \Msg{1.5}{5.6}{-7.0}{\textbf{3} firma de 384 B}
  \Msg{5.6}{9.9}{-7.45}{cuotas de prueba; consume la captura}
  \Paso{5.6}{-7.6}{3.6}{Verifica RSA (perfil Partner), confianza e instante}
  \Paso[pasoError]{9.9}{-7.6}{3.8}{Firma alterada o repetida: rechazo; la captura ya se consumió}
  \draw[msg] (5.6,-8.95) -- node[msgL, pos=0.74]{revalida cuenta, vínculo y confianza} (14.3,-8.95);
  \Msg{5.6}{9.9}{-9.4}{desafío MFA}
  \Msg[msgR]{5.6}{1.5}{-9.85}{desafío, nunca una sesión}
  \Msg{1.5}{5.6}{-10.4}{\textbf{4} código TOTP o de recuperación}
  \Msg{5.6}{9.9}{-10.85}{sesión con procedencia del vínculo}
  \Msg[msgR]{5.6}{1.5}{-11.3}{token de sesión}
  \Paso{1.9}{-11.45}{3.5}{Qadra exige el mismo Owner elegido y con rol Owner}
  \Paso{9.9}{-11.0}{3.8}{Retirar el vínculo o cambiar la confianza la invalida, sin pasarla a contraseña}
\end{tikzpicture}

  \caption{Acceso del Owner por certificado. La firma reemplaza el primer factor y solo concede el desafío MFA del acceso por contraseña; cuenta, vínculo y confianza se revalidan antes del desafío y de la sesión.}
  \label{fig:impl-acceso-certificado}
\end{figure}

\paragraph{Servicio y habilitación.} Las rutas del vínculo y del acceso solo admiten JSON estricto y acotado con evidencia pública, comparten admisión y presupuesto con la API y no reintentan escrituras inciertas; la preparación entrega los bytes exactos que se firman, y el servicio no retiene el token ni modifica campos de usuario. El binario no ejecuta migraciones ni publica confianza, y el alta exige la confianza publicada vigente; registrar un vínculo no enrola cuentas ni habilita el acceso por certificado ni la firma individual, que usa ese vínculo (\cref{sec:impl-firma-individual}). El acceso por certificado es opcional y está desactivado por omisión; sus cuotas y tiempos se validan al arrancar (límites en la \cref{sec:pruebas-vinculo-certificado}) y, si la configuración está incompleta, sus rutas no existen y una ruta pública anuncia que no está disponible. El navegador no puede habilitarlo, y habilitarlo no registra titulares, no publica confianza ni entrega claves.

\paragraph{Interfaz.} \emph{Ingresar con certificado} conserva la contraseña si la disponibilidad es negativa, ausente o inválida, y no usa el recibo elegido como prueba de vigencia. Cada intento firma una declaración nueva, uno incierto o vencido exige otra, y el plazo local monotónico no sustituye al del servidor (\cref{sec:impl-qadra-sesion}).

\paragraph{Límites.} Tras restaurar una copia, una firma nueva solo concede MFA nueva. La admisión de inventario no sustituye el verificador global de auditoría ni acredita una restauración completa. La habilitación privada está en la \cref{sec:impl-portal-ca}, y la verificación, con firmas OpenSSL y no del ayudante Rust/WebAssembly, en la \cref{sec:pruebas-vinculo-certificado}.

\Fuentes{ADR 0067 y 0068; \rutarepo{docs/owner-certificate-login-http.md}; \rutarepo{docs/owner-certificate-login-operations.md}; \rutarepo{docs/owner-certificate-login-interface.md}.}


\subsection{Directorio de miembros y revocación durable}\label{sec:impl-miembros}

El directorio realiza el ciclo de vida de CU-02 y los roles y bajas de RF-02 y RF-03 (\cref{tab:cu-02,tab:rf-acceso}), aparte del directorio de personas de un expediente. El diseño, con la guarda del último Owner, está en la \cref{fig:cu-seq-alta}, y la revocación por generación, en la \cref{fig:impl-identidad-estado}. El alta directa por un Owner (\texttt{POST users}, \cref{tab:impl-http-identidad}) no es una invitación (\cref{sec:impl-invitaciones}): entrega el material MFA una sola vez y, si se pierde la respuesta, no se recupera.

\paragraph{Directorio.} Solo un Owner activo consulta el directorio y los selectores de asignados o disponibles, que muestran identificador, correo, rol, actividad y revisión, sin secretos ni generación. Los filtros preceden a una paginación de 1 a 100 filas por UUID, con cursor ligado a filtros y expediente y sin instantánea común entre páginas.

\paragraph{Cambios.} Rol y actividad cambian juntos contra la revisión esperada: un cambio real incrementa revisión y generación con su auditoría en una transacción, repetir valores vigentes no cambia nada y una revisión anterior da conflicto aunque coincidan. PostgreSQL impide dejar al despacho sin Owner activo, también ante cambios concurrentes. Desactivar conserva asignaciones y autoría; retirar una asignación es otra operación, permitida con el expediente cerrado. Reactivar exige contraseña y MFA, y un resultado incierto, volver a consultar la cuenta (\cref{sec:impl-qadra-conciliacion}).

\paragraph{Límites.} Una petición ya autorizada puede terminar tras la revocación, que no recupera bytes entregados. Restaurar PostgreSQL puede retroceder la generación: antes de reabrir se invalidan sesiones y desafíos del Redis correspondiente (\cref{anexo:miembros}). Ninguna primitiva coordina la restauración conjunta de PostgreSQL y Redis, descrita en la \cref{sec:impl-despliegue}, que los controladores de mantenimiento no automatizan del todo. La verificación está en la \cref{sec:pruebas-miembros}.

\Fuentes{\rutarepo{docs/members-api.md}; \rutarepo{docs/database-operations.md}; \rutarepo{docs/deployment-backups.md}; ADR 0043.}

